\documentclass[11pt,letterpaper]{letter}
\usepackage[margin=1in]{geometry}
\usepackage{hyperref}

\begin{document}

\begin{letter}{Editorial Board \\ \textit{The Fibonacci Quarterly}}

\opening{Dear Editor,}

I am pleased to submit the amended version of ``A Phase-Gated AR(2) Framework for Colonic Crypt Renewal: A Temporal Analogue to Fibonacci Tissue Codes with Tuft Cell Readouts'' for consideration as a Reply article in \textit{The Fibonacci Quarterly}.

This paper proposes a temporal complement to Boman et al.'s Fibonacci tissue code for colonic crypt renewal, using Phase-Gated Autoregressive models of order 2 (PAR(2)) whose companion matrix eigenvalues encode persistence, oscillation, and stability of crypt-level observables. The key contributions are:

\begin{enumerate}
\item A layered molecular kernel architecture (BMAL1/CLOCK phase source, NOTCH fast kernel, WNT/AXIN2 slow kernel) that maps signalling pathways to PAR(2) coefficients.
\item Tuft cells proposed as delayed, accumulative readouts of temporal coherence, with the bidirectional pattern (suppressed in tumours, upregulated post-chemotherapy) explained through eigenvalue hierarchy dynamics.
\item Systematic cross-validation against 22 datasets across 5 species on the PAR(2) Discovery Engine platform, with transparent reporting of confirmed, partially disconfirmed, and untested predictions.
\item Honest treatment of the Fibonacci/golden-ratio connection: individual genes cluster near $1/\varphi = 0.618$ but genome-wide enrichment is not significant ($p = 0.15$).
\end{enumerate}

This amended version incorporates findings from the comprehensive Nguyen, Lausten and Boman (2025) review in \textit{Cells}, corrects the BMAL1-KO prediction based on organoid perturbation data, and addresses the M-cell terminology distinction between Boman's renewal model and the cell-biological literature.

All analyses are reproducible via the publicly available PAR(2) Discovery Engine platform (\url{https://par2-discovery-engine.replit.app}) and archived code on Zenodo (DOI: 10.5281/zenodo.13746927).

The author declares a pending patent on the mathematical framework (Application No.\ GB2518973.9) but no other competing interests. The manuscript has not been submitted elsewhere.

\closing{Sincerely,\\Michael Whiteside\\Independent Researcher, United Kingdom\\mickwh@msn.com}

\end{letter}
\end{document}
